% Irodov, Arihant PDF pp. 35--36, printed pp. 26--27.
\exproblem{151}{26}
图1.39中，两个物块质量为 $m_1,m_2$，以劲度系数
$\kappa$ 的轻弹簧相连，置于光滑水平面上。
物块1起初靠着墙，物块2将弹簧向左压缩一小段
$x$ 后由静止释放。求物块1刚离墙后系统质心的速度。
\exsolution
在弹簧尚未恢复原长之前，弹簧对物块1的压力向左，
墙的支持力使物块1始终静止。物块2单独作弹簧振动，
弹簧恢复自然长度瞬间，物块1刚要离开墙。
在此前阶段，初始弹性能
$\tfrac12\kappa x^2$ 全部变成物块2的动能：
\[
\frac12\kappa x^2=\frac12m_2v_2^2
\quad\Longrightarrow\quad
v_2=x\sqrt{\frac{\kappa}{m_2}}.
\]
此瞬间 $v_1=0$，系统质心速度为
\[
\boxed{v_C=\frac{m_2v_2}{m_1+m_2}
=\frac{x\sqrt{\kappa m_2}}{m_1+m_2}},
\]
方向远离墙。
\excheck 不能从 $t=0$ 起应用\emph{两物块}的总水平
动量守恒，因为在离墙前墙对物块1有外力。

\exproblem{152}{26}
两个物块以自然长 $\ell_0$ 的轻弹簧连接，
弹簧劲度系数 $\kappa$，放在光滑水平面上。
从 $t=0$ 起对右侧物块 $m_2$ 施加恒定水平向右力 $F$。
求两物块距离的最大值与最小值：
（a）当两物块质量相等；
（b）任意质量 $m_1,m_2$。
\exsolution
设两块坐标 $x_1,x_2$，相对距离
$\ell=x_2-x_1$，则
\[
m_1\ddot x_1=\kappa(\ell-\ell_0),\qquad
m_2\ddot x_2=F-\kappa(\ell-\ell_0).
\]
相减得到
\[
\ddot\ell+\kappa\left(\frac1{m_1}+\frac1{m_2}\right)
(\ell-\ell_0)=\frac F{m_2}.
\]
令 $\mu=m_1m_2/(m_1+m_2)$，
$\Omega^2=\kappa/\mu$，
$\delta=Fm_1/[\kappa(m_1+m_2)]$。
用初值 $\ell(0)=\ell_0,\ \dot\ell(0)=0$ 解得
\[
\ell(t)=\ell_0+\delta(1-\cos\Omega t).
\]
因此
\[
\boxed{\ell_{\min}=\ell_0,\qquad
\ell_{\max}=\ell_0+\frac{2m_1F}{\kappa(m_1+m_2)}}.
\]
当 $m_1=m_2$ 时，
$\boxed{\ell_{\max}=\ell_0+F/\kappa}$。
\excheck 任意时间 $1-\cos\Omega t\in[0,2]$，
所以给出的最大、最小距离均可以实际达到。

\exproblem{153}{26}
图1.41中两个质量均为 $m$ 的小立方体上下排列，
以劲度系数 $\kappa$ 的竖直轻弹簧相连，
弹簧初始压缩 $\Delta\ell$，两物块又用绳固定。
烧断绳后，求：（a）下方立方体最终会离地的条件；
（b）若 $\Delta\ell=7mg/\kappa$，
系统质心能上升的最大高度。
\exsolution
定义 $q=mg/\kappa$，取弹簧形变
$x$（拉伸为正）。下方立方体未离地时，上方物块
从 $x=-\Delta\ell$ 由静止运动。
其机械能守恒：
\[
\tfrac12mv^2+\tfrac12\kappa x^2+mgx
=\tfrac12\kappa(\Delta\ell)^2-mg\Delta\ell.
\]
若下块一直贴地，上块最高时 $v=0$，
非平凡转折形变为
$x_{\max}=\Delta\ell-2q$。
下块受弹簧上拉力 $\kappa x$，
只有当 $x>q$ 才可能离地，因此
\[
\boxed{\Delta\ell>3q=\frac{3mg}{\kappa}}.
\]
在刚离地时 $x=q$，上块动能由上述能量式给出
\[
v^2=\frac{\kappa}{m}(\Delta\ell+q)(\Delta\ell-3q).
\]
此刻下块速度仍为零，因此系统质心速度为 $v/2$。
从初始到下块刚离地，质心已升高
$(\Delta\ell+q)/2$；之后系统只受重力，
质心再升高 $v^2/(8g)$，故
\[
\boxed{h_C^{\max}=
\frac{(\Delta\ell+q)^2}{8q}
=\frac{mg}{8\kappa}
\left(1+\frac{\kappa\Delta\ell}{mg}\right)^2}.
\]
代入 $\Delta\ell=7q$：
$\boxed{h_C^{\max}=8mg/\kappa}$。
\excheck 判断下块离地必须比较弹簧\emph{最大拉伸}
而非初始压缩时的弹簧力。
质心在两块离地以后受的加速度始终是 $-\bm g$。

\exproblem{154}{26}
两辆相同小车1、2在平行无摩擦轨道上相向运动，
每辆空车质量为 $M$，各载一名质量 $m$ 的人。
小车相遇时两人沿垂直轨道方向跳到对方车上。
最终车1静止，车2保持原来的运动方向、
速度为 $v$。求两车原来沿轨道的有符号速度 $v_1,v_2$。
\exsolution
跳跃沿横向进行，对运动方向的碰撞可以看作
“外来乘客与车体”之间的非弹性结合，
各乘客在跳出时保留原车的纵向速度。
故最终纵向速度为
\[
V_1=\frac{Mv_1+mv_2}{M+m}=0,\quad
V_2=\frac{Mv_2+mv_1}{M+m}=v.
\]
第一式给出 $v_1=-(m/M)v_2$，
代入第二式后得
\[
\boxed{v_1=-\frac{mv}{M-m},\qquad
v_2=\frac{Mv}{M-m}}.
\]
\excheck 要让车2继续沿初始方向运动且该速度为正，
必须有 $M>m$。两车最初相向：
$v_1$ 与 $v_2$ 确为相反符号。

\exproblem{155}{27}
两辆相同空车质量为 $M$，
沿同一直线无摩擦轨道同向以速率 $v_0$ 行驶。
后车载一名质量 $m$ 的乘客，乘客以相对后车的
向前速度 $u$ 跳到前车。
求两车最终速度。
\exsolution
设乘客起跳后后车速度为 $V_r$。
乘客离后车时，其纵向速度为 $V_r+u$。
对起跳过程的后车—乘客体系守恒纵向动量：
\[
(M+m)v_0=MV_r+m(V_r+u),
\]
故
\[
\boxed{V_r=v_0-\frac{m}{M+m}u}.
\]
乘客随后落到前车上，为完全非弹性纵向结合，
落地前乘客速度为
$v_p=V_r+u=v_0+Mu/(M+m)$。
对前车及乘客守恒纵向动量：
\[
(M+m)V_f=Mv_0+mv_p,
\]
所以
\[
\boxed{V_f=v_0+\frac{mM}{(M+m)^2}u}.
\]
\excheck 后车减速、前车加速；整体纵向动量
在两次相互作用全过程仍然守恒。

\exproblem{156}{27}
一辆质量为 $M$ 的无摩擦小车上静止站着两个人，
每人质量均为 $m$。
两人分别以相对小车的水平速度 $u$
向同一方向跳下。求小车最终速度：
（1）两人同时跳；
（2）两人先后跳；
并比较两种情况下的小车速度大小。
\exsolution
取乘客跳出方向为正。
（1）同时跳出时，小车末速 $V_{\rm sim}$，
两人末速均为 $V_{\rm sim}+u$，动量守恒：
\[
0=MV_{\rm sim}+2m(V_{\rm sim}+u),
\quad
\boxed{V_{\rm sim}=-\frac{2m}{M+2m}u}.
\]
（2）先跳者离车时，仍留在车上的人
与空车总质量 $M+m$，末速为 $V_1$。
第一跳：
$0=(M+m)V_1+m(V_1+u)$，
故 $V_1=-mu/(M+2m)$。
第二跳时，由相对于第一跳前的运动系守恒动量，
空车最终速度再减少 $mu/(M+m)$，于是
\[
\boxed{V_{\rm seq}=
-\frac{m(2M+3m)}{(M+m)(M+2m)}u}.
\]
二者速度大小之比为
\[
\boxed{\frac{|V_{\rm seq}|}{|V_{\rm sim}|}
=1+\frac{m}{2(M+m)}>1}.
\]
\excheck 先后跳下时第二人面对的可反冲剩余总质量更小，
故同样相对跳速产生更大最终反冲。

\exproblem{157}{27}
一条均匀柔软链条竖直悬挂，下端恰触到桌面。
烧断上端支承绳后，链条开始自由落下并逐渐堆积在桌上。
证明在任意时刻，\emph{仍在下落的链段}对桌面的冲击力
是\emph{已经静置在桌上链段}重力的两倍。
\exsolution
设链条线密度为 $\lambda$，
在时刻 $t$ 已经到达桌面的链长为 $x$。
由于烧断后各链元由静止开始自由落体，
达到桌面的链元速率为 $v=\sqrt{2gx}$。
单位时间内到达桌面的链条质量
$\dot m=\lambda v$。
下落链元被桌面制止，需要动量变化率
\[
F_{\rm impact}=\dot m\,v=\lambda v^2=2\lambda gx.
\]
已经堆积在桌面上的链条质量为 $\lambda x$，
其静重力为
$W_{\rm rested}=\lambda gx$，因此
\[
\boxed{F_{\rm impact}=2W_{\rm rested}}.
\]
\excheck 全部对桌面的瞬时压力实际上为
$F_{\rm impact}+W_{\rm rested}=3\lambda gx$，
但题目只把\emph{下落部分的冲击力}与静置部分重量相比。
假定链元落桌后立即静止且无弹跳。

\exproblem{158}{27}
一颗钢球质量 $m=50\,\mathrm g$，
由高度 $h=1.0\,\mathrm m$ 落向很重的水平板。
钢球每次碰撞后的反弹速率
是碰撞前速率的 $1/\eta$，
其中 $\eta=1.25$。
求经无数次弹跳，钢球传给平板的总冲量。
\exsolution
第一次碰撞前的下降速率
$v_0=\sqrt{2gh}$。
第 $n$ 次（从 $n=0$ 编号）碰撞前速率
$v_n=v_0/\eta^n$，碰后向上速率 $v_n/\eta$。
按向下为正，钢球碰撞动量变化为
$-m v_n/\eta-(+m v_n)$，
所以钢球传给平板的向下冲量大小为
$m v_n(1+1/\eta)$。
总冲量是收敛几何级数：
\[
J=\sum_{n=0}^\infty m v_0
\left(1+\frac1\eta\right)\eta^{-n}
=mv_0\frac{\eta+1}{\eta-1}.
\]
代入数值：
\[
\boxed{J=
0.050\sqrt{2(9.81)(1.0)}
\frac{2.25}{0.25}
\simeq1.99\,\mathrm{N\,s}},
\]
方向向下。
\excheck 第一次碰撞的冲量约 $0.399\,\mathrm{N\,s}$，
后续总和约 $1.60\,\mathrm{N\,s}$，符合几何级数。
原书印刷答案页确实写有 $0.2\,\mathrm{kg\,m/s}$，但将题面参数代入同页公式应得约 $2.0\,\mathrm{kg\,m/s}$；因此这是印刷数值勘误，并非 PDF 文字提取错误。

\exproblem{159}{27}
质量为 $M$ 的木筏与质量为 $m$ 的人一起静止浮在湖面，
忽略水的水平阻力。
人在筏上相对木筏移动总距离 $\ell'$，
相对速度随时间变化为 $v'(t)$，最后停住。
求：（a）木筏相对湖岸的位移；
（b）人对木筏的水平作用力。
\exsolution
取人的移动方向为正，筏岸位移记 $X$，
人岸位移为 $X+\ell'(t)$。
系统初始总水平动量为零，始终没有外部水平合力：
\[
M\dot X+m(\dot X+v')=0
\Longrightarrow
\dot X=-\frac{m}{M+m}v'.
\]
积分得木筏的最终位移
\[
\boxed{\Delta X=-\frac{m}{M+m}\ell'}.
\]
对木筏单独应用牛顿第二定律，
人施加的水平作用力为
\[
\boxed{F_{\rm man\to raft}=M\ddot X
=-\frac{mM}{M+m}\frac{\mathrm dv'}{\mathrm dt}}.
\]
\excheck 人在木筏上加速前行时筏向后受力，
人在末段减速时木筏可能受向前的力；
人停下后系统质心仍保持原水平位置。

\exproblem{160}{27}
一固定滑轮上跨轻绳，一端悬挂梯子和攀登者，
另一端悬挂质量为 $M$ 的配重。
攀登者质量为 $m$，在梯子上相对梯子向上爬
$\ell'$ 后停住，忽略滑轮摩擦与绳质量。
求这个系统质心的竖直位移。
\exsolution
为了使题目无需给梯子的质量还能得到唯一结果，
应采用原书常见的\emph{初始平衡配重}模型：
梯子与人的总质量等于配重 $M$，
即梯子质量为 $M-m$。
设梯子相对地面竖直位移为 $y$（向上为正），
则另一端配重因绳长约束位移为 $-y$；
人相对梯子额外上升 $\ell'$，
故其总位移为 $y+\ell'$。
总质量为 $2M$，于是质心位移
\[
\Delta y_C=
\frac{(M-m)y+m(y+\ell')+M(-y)}{2M}
=\boxed{\frac{m\ell'}{2M}}.
\]
\excheck 梯子自身的未知位移 $y$ 在质心加权和中相消。
若原题并未隐含梯子与人总质量等于配重，
则单给 $M,m,\ell'$ 不足以确定质心位移；
本解明确补出该配重假设，待独立版本对校。
